\chapter{Solución de Problemas}
\label{ch:troubleshooting}

Si tu proyecto no funciona como se espera, este capítulo te ayuda a identificar problemas comunes en
\textbf{ambos lados del sistema}:

\begin{itemize}
    \item \textbf{Android (EmitterApp):} controles de la interfaz, emparejamiento/conexión Bluetooth y si se están procesando comandos/telemetría.
    \item \textbf{ESP32 (ESP32\_BT\_Controller):} carga del firmware, cableado de pines, integridad de alimentación, enlace Bluetooth SPP y generación de telemetría.
\end{itemize}

\section{Antes de Empezar: Lista Rápida}
\label{sec:quick-checklist}

\begin{itemize}
    \item Confirma que el ESP32 está alimentado y estable (sin brownouts/resets).
    \item Confirma que tu dispositivo Android está emparejado con el nombre Bluetooth correcto (por defecto es \texttt{ESP32-BT}).
    \item Mantén el teléfono a pocos metros para las pruebas iniciales.
    \item Asegúrate de que \textbf{todas las tierras sean comunes} (ESP32, driver de motor, sensores, fuente externa de 5V/alimentación de servos).
    \item Si estás usando pines sensibles al arranque (por ejemplo, GPIO0), desconecta temporalmente los circuitos externos durante el flasheo.
\end{itemize}

\section{El Firmware del ESP32 No Se Carga}
\label{sec:esp32-upload}

Si el firmware no se carga correctamente:

\begin{itemize}
    \item Verifica que la placa/target correcto esté seleccionado en tu IDE/herramientas.
    \item Confirma que el puerto serie correcto esté seleccionado.
    \item Prueba con otro cable USB (muchos cables ``solo carga'' causan fallos de carga).
    \item Mantén presionado \textbf{BOOT} durante la carga si tu placa requiere entrada manual al bootloader.
    \item Desconecta el cableado externo de los \textbf{pines de arranque (boot-strap)} durante el flasheo (los más comunes son GPIO0 y GPIO2).
\end{itemize}

\subsection*{Notas sobre pines de arranque (importante)}
La configuración de pines del firmware usa GPIO0 como entrada digital de indicador (\texttt{PIN\_IND4 = GPIO0}) en
\texttt{FULL\_RC\_MODE}. GPIO0 es un pin de arranque: resistencias pull-up/pull-down externas pueden impedir el arranque o el flasheo.
Si el flasheo falla, desconecta cualquier cosa conectada a GPIO0, flashea de nuevo y luego vuelve a conectar.

\section{El ESP32 Enciende pero Se Reinicia Aleatoriamente}
\label{sec:esp32-resets}

Los reinicios aleatorios suelen estar relacionados con la alimentación:

\begin{itemize}
    \item Los motores/servos pueden bajar el voltaje de forma repentina. Usa un regulador dedicado para el ESP32 y una fuente separada para motores/servos.
    \item Asegúrate de que la fuente del ESP32 pueda entregar corriente pico (ráfagas de Wi-Fi/Bluetooth pueden causar caídas).
    \item Agrega capacitancia de reserva cerca de drivers de motor/riel de servos (p.\ ej., 470--1000~\(\mu\)F) y un capacitor cerámico más pequeño cerca del ESP32.
    \item Revisa si los reguladores de las placas de desarrollo se sobrecalientan cuando se alimentan servos desde el mismo riel de 5V (no recomendado).
\end{itemize}

\section{Problemas de Emparejamiento y Conexión Bluetooth (Android \texorpdfstring{$\leftrightarrow$}{<->} ESP32)}
\label{sec:bluetooth-problems}

\subsection*{Síntoma: El dispositivo no aparece en el escaneo}
\begin{itemize}
    \item Confirma que el ESP32 está alimentado y ejecutándose (LED/indicadores de actividad si están disponibles).
    \item Asegúrate de que el teléfono esté cerca del ESP32.
    \item Apaga y vuelve a encender el ESP32.
\end{itemize}

\subsection*{Síntoma: El emparejamiento pide PIN o falla}
El PIN por defecto normalmente es \texttt{1234}.
\begin{itemize}
    \item Elimina (``olvida'') el dispositivo en la configuración Bluetooth de Android y empareja de nuevo.
    \item Intenta emparejar primero en la configuración del sistema Android y luego conéctate desde la app.
\end{itemize}

\subsection*{Síntoma: La app dice ``Connected'' pero los controles no hacen nada}
\begin{itemize}
    \item Verifica que te conectaste al dispositivo correcto (si hay varios \texttt{ESP32-BT} cerca).
    \item Confirma que el ESP32 está ejecutando el build de firmware esperado (no un sketch viejo).
    \item Del lado del ESP32, usa logging serial (si está habilitado) para confirmar que llegan paquetes de comandos.
\end{itemize}

\subsection*{Síntoma: Se conecta y luego se desconecta repetidamente}
\begin{itemize}
    \item Esto suele ser un problema de \textbf{estabilidad de alimentación} en el ESP32 (ver Sección~\ref{sec:esp32-resets}).
    \item Reduce la carga de motores/servos durante las pruebas iniciales (desconecta motores, prueba primero con LEDs).
    \item Mantén el teléfono cerca y evita carcasas metálicas que bloqueen Bluetooth.
\end{itemize}

\section{Los Controles Se Sienten Mal en la App (Lado Android)}
\label{sec:android-controls}

\subsection*{El joystick vuelve al centro vs se queda fijo}
EmitterApp soporta diferentes modos de joystick (por ejemplo, spring vs hold). Si el robot se comporta
como si la aceleración ``no se quedara'' o la dirección estuviera ``pegada'':

\begin{itemize}
    \item Revisa los modos de stick seleccionados en la configuración de la app.
    \item Para aceleración, muchos usuarios prefieren el comportamiento \textbf{hold}; para dirección, muchos prefieren \textbf{spring}.
\end{itemize}

\subsection*{Los botones/interruptores aparecen pero no afectan el hardware}
\begin{itemize}
    \item Confirma que el firmware del ESP32 mapea esos canales a pines/salidas reales en el modo de proyecto elegido.
    \item Si usas \texttt{FULL\_RC\_MODE\_MCP}, verifica el cableado del MCP23017 y la configuración de la dirección I2C (dirección por defecto \texttt{0x20}).
\end{itemize}

\section{Los Motores No Se Mueven}
\label{sec:motors}

\subsection*{Verificaciones eléctricas}
\begin{itemize}
    \item Revisa la alimentación del driver de motor y la polaridad.
    \item Confirma que la \textbf{tierra lógica} del driver esté conectada a la tierra del ESP32.
    \item Confirma que las salidas del driver estén conectadas a los terminales del motor (invierte cables si la dirección está al revés).
\end{itemize}

\subsection*{Verificación de señales}
\begin{itemize}
    \item Confirma que los pines PWM coincidan con tu cableado. En el mapeo de pines por defecto, los canales PWM incluyen:
    GPIO18/19/23/25/26/27.
    \item Confirma que los pines de dirección estén conectados correctamente:
    \texttt{PIN\_MD\_LEFT = GPIO12} y \texttt{PIN\_MD\_RIGHT = GPIO15} (ambos deben mantenerse en low al arranque).
    \item Si tu driver necesita dos pines de dirección por motor (IN1/IN2), debes adaptar el cableado/el firmware en consecuencia.
\end{itemize}

\subsection*{Consejo práctico de depuración}
Prueba las salidas de motor con un simple ``LED + resistencia'' en el pin PWM primero (deberías ver cambios de brillo con el movimiento del stick),
y luego conecta el driver de motor.

\section{El Servo No Responde o Tiembla}
\label{sec:servo}

\begin{itemize}
    \item Asegúrate de que el servo tenga una fuente de 5V adecuada ( \textbf{no} alimentes servos medianos/grandes desde la placa del ESP32).
    \item Confirma que la tierra del servo esté unida a la tierra del ESP32.
    \item Verifica que la señal de control esté cableada a un pin de salida PWM configurado (por ejemplo, GPIO23 o GPIO25).
    \item Si el servo tiembla, agrega una mejor fuente de alimentación y reduce el ruido eléctrico (cables más cortos, capacitores, rieles de alimentación separados).
\end{itemize}

\section{Los Datos de Telemetría No Aparecen en la App}
\label{sec:telemetry}

\subsection*{Síntoma: Conectado, pero los widgets de telemetría no se actualizan}
\begin{itemize}
    \item Espera unos segundos después de conectar: el ESP32 envía la configuración a la app en la primera conexión.
    \item Verifica que el sensor esté cableado al tipo de entrada correcto:
    \begin{itemize}
        \item Telemetría analógica: usa pines ADC1 (p.\ ej., GPIO34/35/36/39 en \texttt{FULL\_RC\_MODE}).
        \item Indicadores digitales: GPIO4/5/2/0 (nota: precaución de arranque en GPIO0).
        \item Sensores I2C: SDA GPIO21, SCL GPIO22.
    \end{itemize}
    \item Si estás usando \texttt{FULL\_RC\_MODE\_MCP}, confirma que el MCP23017 se detecte y esté cableado correctamente.
\end{itemize}

\subsection*{Síntoma: La telemetría está presente pero los valores se ven mal}
\begin{itemize}
    \item Revisa el escalado del sensor (divisores de voltaje, rango del ADC, unidades).
    \item Reduce el ruido: acorta cables analógicos, agrega filtrado (filtro RC) y evita tender líneas analógicas junto a cables de motores.
\end{itemize}

\section{Dispositivos I2C No Detectados (MCP23017 o Sensores)}
\label{sec:i2c}

\begin{itemize}
    \item Verifica que SDA esté en GPIO21 y SCL en GPIO22 (según \texttt{pins.h}).
    \item Confirma que el dispositivo use lógica de 3.3V (o usa conversión de nivel).
    \item Revisa resistencias pull-up en SDA/SCL (muchas placas incluyen; demasiadas pull-ups en paralelo también pueden causar problemas).
    \item Confirma la dirección I2C (por defecto del MCP23017 en el repositorio es \texttt{0x20}).
\end{itemize}

\section{Consejos Generales de Depuración (Flujo Recomendado)}
\label{sec:debugging-tips}

\begin{enumerate}
    \item \textbf{Empieza por el enlace:} verifica primero el emparejamiento Bluetooth y una conexión estable.
    \item \textbf{Prueba salidas sin motores:} LEDs en pines PWM/de switch confirman que la ruta de control funciona.
    \item \textbf{Agrega un subsistema a la vez:} motores, luego servos, luego sensores, luego expansión I2C.
    \item \textbf{Usa logging:}
    \begin{itemize}
        \item ESP32: salida debug por serial (cuando está habilitada en el firmware).
        \item Android: revisa los indicadores para confirmar el estado de conexión y la telemetría entrante.
    \end{itemize}
    \item \textbf{Trata la alimentación como un componente de primera clase:} la mayoría de fallos intermitentes provienen de subvoltaje o fuentes ruidosas.
\end{enumerate}